\title{Kinematic Maps of the LMC's Stellar Halo with Spec-S5}

\author{Hayden R. Foote}
\affil{Steward Observatory, University of Arizona}

\author{Himansh Rathore}
\affil{Steward Observatory, University of Arizona}

\author{Gurtina Besla}
\affil{Steward Observatory, University of Arizona}

\textbf{Facts:} The LMC-SMC system is the nearest example of an interacting galaxy pair, allowing us to study the process of heirarchichal galaxy assembly in our cosmic backyard. A galaxy's stellar halo is a key record of its assembly history, as it is composed primarily of accreted stellar populations, and its kinematics encodes perturbations from ongoing interactions/mergers and the distribution of its host's dark matter.

\textbf{Problem:} While several low-surface brightness stellar populations consistent with a stellar halo around the LMC have been identified in e.g., SMASH (Nidever et al. 2019), kinematic data is required to unambiguously identify these structures as the LMC stellar halo as opposed to tidal debris from the LMC-SMC interaction. The dearth of existing kinemtic data for the diffuse outer regions of the LMC is driven primarily by faintness -- while SMASH reaches a \textit{g}-band depth of 24.8, Gaia and existing state-of-the-art spectroscopic surveys like DESI cannot provide kinematics for targets with \textit{g}$\gtrsim 20$. 

\textbf{Solution:} Next-generation astrometric and spectroscopic surveys are required to provide 6-D phase-space information on the diffuse outer LMC. In particular, \textit{Roman} proper motions coupled with Spec-S5 line-of-sight velocities will measure precision kinematics at greater depths [$g\sim22$ gets to the tip of the main sequence, need to check if Spec-S5 can do this], providing a comprehensive kinematic characterization of the LMC's outskirts for the first time. 

\textbf{Strategy, Forecasts, and Implications}: Here, we use our group's library of constrained N-body simulations of the LMC-SMC-MW interaction to identify potentially observable kinematic signatures in the LMC's halo and the influence of the MW and SMC. With precise 6-D phase-space measurements of the LMC stellar halo from \textit{Roman} and Spec-S5, we aim to accomplish the following goals.

\textit{Kinematic Identification of the LMC Stellar Halo}: We begin in \texttt{kinematic_halo_ids.ipynb}, in which we analyze the present-day snapshot of the \texttt{lmcsmcmw09} MEGHA simulation (Rathore et al. in prep.). To prove that we can easily separate LMC halo stars from MW halo stars kinematically, we plot LMC and MW halo particles within 15 kpc of the LMC's center of mass in various phase-space projections. We find that $E vs. L_x$ and $v_y vs. v_z$ spaces are effective at kinematically separating MW and LMC particles. Additionally, we plot SMC stellar debris (SMC stars more than 5 kpc from the SMC's center of mass), and find they are kinematically similar (in a Galactocentric frame) to LMC halo particles, though the SMC's tidal tails form coherent structures in phase-space.

$\Lambda$CDM cosmology predicts that LMC-mass galaxies should host stellar halos (e.g., Tau et al. 2024). Conclusively identifying (or failing to identify) a diffuse stellar halo around the LMC is a direct test of this prediction, and also opens many opportunities to constrain the assembly history of the LMC and MW. While the LMC is generally understood to be on its first infall to the MW (e.g., Besla et al. 2007), a second-passage orbit remains possible (Vasiliev 2024). If the LMC is on a second-passage orbit, its stellar halo would likely have been stripped during its first pericenter, so the presence of a stellar halo around the LMC would strongly suggest a first-infall orbit.

\textit{Identification of Kinemtic Signtures of Milky Way Tides}: After proving the LMC's stellar halo should be identifiable, we plot the velocity signatures of the MW's tides and SMC perturbations in \texttt{SMC_perturbations.ipynb}. Here, we select LMC halo particles in thin spherical shells at different distances from the LMC center, plotting their radial velocity and tangential dispersion fields in Mollweide projection, identifying a quadrupolar radial velocity signature of >20 km/s at distances >30 kpc due to the Milky Way tides.

The quadrupolar radial velocity signature of the Milky Way's tidal influence will constrain the MW/LMC mass ratio and the LMC's orbit (e.g., Garavito-Camargo et al. 2021).  

\textit{Identifying the LMC Halo's Response to the SMC}: We compare the maps of the LMC halo's velocity quadrupole to similar plots made with a simulation that does not include the SMC to isolate the effects of Milky Way tides. We find that the SMC leaves identifiable modifications to the signature of the MW's tidal field in our simulations. Remaining work-in-progress includes adding forecasted Spec-S5 errors to determine whether these features remain identifiable with a realistic error budget. 

The SMC is expected to perturb the LMC's dark matter and stellar halo during their interaction (Foote et al. 2026). Identification of these perturbations to the MW's tidal signature from the SMC will confirm the SMC's influence on the LMC dark matter and stellar halos, opening the door for subsequent studies using the LMC-SMC interaction to constrain DM physics. 